\chapter{Introduction to Mathematical Language}

Math 135 is, for most people, the first course where the \emph{proof} matters more than the
answer. Before we can prove anything we need a precise language: what a statement is, how
statements combine, and how quantifiers such as ``for all'' and ``there exists'' behave.
This chapter sets up that vocabulary.

\section{Statements}

\begin{definition}[Statement]
A \emph{statement} (or proposition) is a sentence that is either true or false, but not both.
\end{definition}

\begin{example}
\begin{itemize}
  \item ``$7$ is prime'' is a true statement.
  \item ``$2^{10} < 1000$'' is a false statement ($2^{10} = 1024$).
  \item ``Every even integer greater than $2$ is a sum of two primes'' (Goldbach) is a
        statement; nobody knows its truth value, but it has one.
  \item ``$x^2 = 4$'' is \emph{not} a statement: its truth depends on $x$. It is an
        \emph{open sentence}, written $P(x)$.
  \item ``Prove the theorem.'' is not a statement (it is a command).
\end{itemize}
\end{example}

An open sentence $P(x)$ becomes a statement once we substitute a value for $x$ or once we
\emph{quantify} $x$ (see below). The set of values $x$ is allowed to range over is the
\emph{domain} (or universe of discourse).

\section{Compound statements and truth tables}

Given statements $A$ and $B$ we form new statements with the logical connectives:

\begin{center}
\begin{tabular}{lll}
\toprule
Symbol & Name & Read as \\
\midrule
$\neg A$ & negation & ``not $A$'' \\
$A \wedge B$ & conjunction & ``$A$ and $B$'' \\
$A \vee B$ & disjunction & ``$A$ or $B$'' (inclusive) \\
$A \Rightarrow B$ & implication & ``if $A$ then $B$'' \\
$A \Leftrightarrow B$ & biconditional & ``$A$ if and only if $B$'' \\
\bottomrule
\end{tabular}
\end{center}

Their meanings are fixed by the truth table:

\begin{center}
\begin{tabular}{cc|ccccc}
$A$ & $B$ & $\neg A$ & $A\wedge B$ & $A \vee B$ & $A\Rightarrow B$ & $A\Leftrightarrow B$\\
\hline
$\TT$ & $\TT$ & $\FF$ & $\TT$ & $\TT$ & $\TT$ & $\TT$ \\
$\TT$ & $\FF$ & $\FF$ & $\FF$ & $\TT$ & $\FF$ & $\FF$ \\
$\FF$ & $\TT$ & $\TT$ & $\FF$ & $\TT$ & $\TT$ & $\FF$ \\
$\FF$ & $\FF$ & $\TT$ & $\FF$ & $\FF$ & $\TT$ & $\TT$ \\
\end{tabular}
\end{center}

\begin{important}[The implication is only false in one row]
$A \Rightarrow B$ is false \emph{only} when $A$ is true and $B$ is false. In particular,
an implication with a false hypothesis is \emph{vacuously true}: ``if $1 = 2$ then I am the
Pope'' is a true statement.
\end{important}

In $A \Rightarrow B$, $A$ is the \emph{hypothesis} and $B$ the \emph{conclusion}. Equivalent
English phrasings: ``$A$ implies $B$'', ``$B$ whenever $A$'', ``$A$ is sufficient for $B$'',
``$B$ is necessary for $A$'', ``$A$ only if $B$''.

\subsection{Logical equivalence}

\begin{definition}
Two compound statements are \emph{logically equivalent}, written $S_1 \equiv S_2$, if they
have the same truth value for every assignment of truth values to their component
statements. A statement that is true in every row is a \emph{tautology}; one that is false
in every row is a \emph{contradiction}.
\end{definition}

\begin{proposition}[Useful equivalences]
For statements $A,B,C$:
\begin{enumerate}
  \item $\neg(\neg A) \equiv A$ (double negation);
  \item $\neg(A\wedge B) \equiv \neg A \vee \neg B$ and $\neg(A \vee B) \equiv \neg A
        \wedge \neg B$ (De Morgan's laws);
  \item $A \wedge (B \vee C) \equiv (A\wedge B)\vee(A\wedge C)$ and
        $A \vee (B \wedge C) \equiv (A\vee B)\wedge(A\vee C)$ (distributive laws);
  \item $A \Rightarrow B \equiv \neg A \vee B$;
  \item $\neg(A \Rightarrow B) \equiv A \wedge \neg B$;
  \item $A \Leftrightarrow B \equiv (A\Rightarrow B)\wedge(B \Rightarrow A)$.
\end{enumerate}
\end{proposition}

\begin{proof}[Proof of (4)]
Compare columns:
\begin{center}
\begin{tabular}{cc|cc|c}
$A$ & $B$ & $\neg A$ & $\neg A \vee B$ & $A \Rightarrow B$ \\
\hline
$\TT$ & $\TT$ & $\FF$ & $\TT$ & $\TT$\\
$\TT$ & $\FF$ & $\FF$ & $\FF$ & $\FF$\\
$\FF$ & $\TT$ & $\TT$ & $\TT$ & $\TT$\\
$\FF$ & $\FF$ & $\TT$ & $\TT$ & $\TT$\\
\end{tabular}
\end{center}
The last two columns agree in every row. The rest are checked the same way.
\end{proof}

\begin{example}
Show $(A \Rightarrow B)\wedge(A \Rightarrow C) \equiv A \Rightarrow (B\wedge C)$ without a
table:
\begin{align*}
(A \Rightarrow B)\wedge(A \Rightarrow C)
  &\equiv (\neg A \vee B)\wedge(\neg A \vee C) \\
  &\equiv \neg A \vee (B \wedge C) &&\text{(distributive law)}\\
  &\equiv A \Rightarrow (B \wedge C).
\end{align*}
\end{example}

\section{Converse and contrapositive}

\begin{definition}
For the implication $A \Rightarrow B$:
\begin{itemize}
  \item the \emph{converse} is $B \Rightarrow A$;
  \item the \emph{contrapositive} is $\neg B \Rightarrow \neg A$;
  \item the \emph{inverse} is $\neg A \Rightarrow \neg B$.
\end{itemize}
\end{definition}

\begin{theorem}
An implication is logically equivalent to its contrapositive:
$A \Rightarrow B \equiv \neg B \Rightarrow \neg A$.
\end{theorem}

\begin{proof}
$\neg B \Rightarrow \neg A \equiv \neg(\neg B) \vee \neg A \equiv B \vee \neg A \equiv
\neg A \vee B \equiv A \Rightarrow B$.
\end{proof}

\begin{important}[Converse is not equivalent]
The converse is \emph{not} equivalent to the original. ``If $n = 4$ then $n$ is even'' is
true; its converse ``if $n$ is even then $n = 4$'' is false (take $n = 2$).
The inverse is the contrapositive of the converse, so the inverse is equivalent to the
converse.
\end{important}

\section{Quantifiers}

\begin{definition}
Let $P(x)$ be an open sentence with domain $S$.
\begin{itemize}
  \item The \emph{universal} statement $\forall x \in S,\ P(x)$ is true if $P(x)$ is true for
        every $x \in S$.
  \item The \emph{existential} statement $\exists x \in S,\ P(x)$ is true if $P(x)$ is true
        for at least one $x \in S$.
\end{itemize}
\end{definition}

\begin{example}
\begin{itemize}
  \item $\forall x \in \R,\ x^2 \geq 0$ is true.
  \item $\exists x \in \R,\ x^2 = -1$ is false; but $\exists x \in \C,\ x^2 = -1$ is true.
        The domain matters.
  \item $\forall n \in \N,\ n^2 + n + 41$ is prime is false: at $n = 40$ we get
        $40^2 + 40 + 41 = 1681 = 41^2$.
\end{itemize}
\end{example}

\subsection{Nested quantifiers}

The order of quantifiers of different types matters.

\begin{example}
\begin{itemize}
  \item $\forall x \in \R,\ \exists y \in \R,\ x + y = 0$ is true: given $x$, choose
        $y = -x$. Here $y$ may depend on $x$.
  \item $\exists y \in \R,\ \forall x \in \R,\ x + y = 0$ is false: a single $y$ would have
        to satisfy $1 + y = 0$ and $2 + y = 0$ at once.
\end{itemize}
Quantifiers of the \emph{same} type can be swapped: $\forall x\, \forall y$ is the same as
$\forall y\, \forall x$.
\end{example}

\subsection{Negating quantified statements}

\begin{theorem}
\[
\neg\bigl(\forall x \in S,\ P(x)\bigr) \equiv \exists x \in S,\ \neg P(x),
\qquad
\neg\bigl(\exists x \in S,\ P(x)\bigr) \equiv \forall x \in S,\ \neg P(x).
\]
\end{theorem}

To negate a long statement, move the negation from left to right, flipping each quantifier
and finally negating the innermost open sentence (using De Morgan's laws and
$\neg(A\Rightarrow B)\equiv A \wedge \neg B$).

\begin{example}
Negate ``for every $\varepsilon > 0$ there is a $\delta > 0$ such that for all $x\in\R$, if
$\abs{x} < \delta$ then $\abs{f(x)} < \varepsilon$''.
\begin{align*}
&\neg\bigl(\forall \varepsilon > 0,\ \exists \delta > 0,\ \forall x \in \R,\
  (\abs{x} < \delta \Rightarrow \abs{f(x)} < \varepsilon)\bigr)\\
&\equiv \exists \varepsilon > 0,\ \forall \delta > 0,\ \exists x \in \R,\
  \neg(\abs{x} < \delta \Rightarrow \abs{f(x)} < \varepsilon)\\
&\equiv \exists \varepsilon > 0,\ \forall \delta > 0,\ \exists x \in \R,\
  (\abs{x} < \delta \wedge \abs{f(x)} \geq \varepsilon).
\end{align*}
\end{example}

\begin{example}
Negate ``every integer greater than $1$ has a prime divisor'': ``there exists an integer
greater than $1$ that has no prime divisor''. Note the negation is \emph{not} ``no integer
greater than $1$ has a prime divisor''.
\end{example}

\section{Reading and writing mathematics}

A theorem typically has the shape ``for all $x$ in $S$, if $H(x)$ then $C(x)$'', even when the
quantifier is hidden: ``if $n$ is odd then $n^2$ is odd'' really means
$\forall n \in \Z,\ (n \text{ odd} \Rightarrow n^2 \text{ odd})$. Before proving something,
identify:
\begin{enumerate}
  \item the quantifiers and the domain;
  \item the hypothesis (what we may assume);
  \item the conclusion (what we must show).
\end{enumerate}

\begin{aside}
Common notation: $\N = \{1,2,3,\dots\}$ (in Math 135 the natural numbers start at $1$),
$\Z$ integers, $\Q$ rationals, $\R$ reals, $\C$ complex numbers. ``Such that'' is often
abbreviated ``s.t.''; $\therefore$ means ``therefore''.
\end{aside}

\section{Common pitfalls}
\begin{itemize}
  \item Treating ``or'' as exclusive. In mathematics $A \vee B$ is true when both are true.
  \item Thinking a false hypothesis makes an implication false. It makes it (vacuously) true.
  \item Proving the converse instead of the statement.
  \item Negating $\forall x,\ P(x)$ as $\forall x,\ \neg P(x)$. The correct negation needs only
        one counterexample.
  \item Swapping $\forall\exists$ to $\exists\forall$, which usually makes a statement much
        stronger.
\end{itemize}

\section{Summary}
\begin{itemize}
  \item A statement is true or false; an open sentence becomes a statement when its variables
        are assigned or quantified.
  \item Connectives $\neg,\wedge,\vee,\Rightarrow,\Leftrightarrow$ are defined by truth tables;
        $A \Rightarrow B$ is false only when $A$ is true and $B$ false.
  \item $A\Rightarrow B \equiv \neg A \vee B \equiv \neg B \Rightarrow \neg A$ (contrapositive);
        the converse $B \Rightarrow A$ is a different statement.
  \item Negation flips quantifiers: $\neg \forall \equiv \exists \neg$ and
        $\neg\exists \equiv \forall \neg$; the order of mixed quantifiers matters.
\end{itemize}
